\section{Experiments}
\label{sec:exp}

\subsection{Setup}
\label{sec:exp:setup}

\begin{table}[t]
\centering
\small
\caption{The runs. \emph{Fail} is the fraction of generations that fail and
\emph{Pos.}\ the number of labelled failures. The attention column gives the
number of layers, the number of query heads, and whether every layer computes
softmax attention or only a subset does.}
\label{tab:runs}
\begin{tabular}{llrrrrl}
\toprule
Model & Data & $N$ & Fail & Pos. & Neg. & Attention \\
\midrule
Llama-3.2-1B & Glaive & 651 & 47\% & 303 & 348 & 16L, 32H dense \\
Llama-3.2-3B & Glaive & 651 & 18\% & 119 & 532 & 28L, 24H dense \\
Gemma-3-1B & Glaive & 651 & 29\% & 190 & 461 & 26L, 4H sliding \\
Llama-3.2-1B & BFCL & 850 & 72\% & 612 & 238 & 16L, 32H dense \\
Llama-3.2-3B & BFCL & 850 & 57\% & 488 & 362 & 28L, 24H dense \\
Qwen3-1.7B & BFCL & 850 & 22\% & 186 & 664 & 28L, 16H dense \\
MiniCPM5-2B & BFCL & 850 & 11\% & 92 & 758 & 42L, 16H dense \\
Qwen3.5-0.8B & BFCL & 850 & 39\% & 328 & 522 & 24L, 6 attentive \\
Qwen3.5-4B & BFCL & 850 & 14\% & 116 & 734 & 32L, 8 attentive \\
Qwen3.5-2B$^\dagger$ & Glaive & 740 & 3\% & 21 & 719 & 24L, 6 attentive \\
Llama-3.2-1B & BFCL-live & 765 & 73\% & 559 & 206 & 16L, 32H dense \\
MiniCPM5-2B & BFCL-live & 779 & 23\% & 178 & 601 & 42L, 16H dense \\
\bottomrule
\end{tabular}

\end{table}

Table~\ref{tab:runs} lists the runs. They cover five model families between
0.8B and 4B parameters and three data distributions: the Glaive
function-calling corpus \citep{glaive2024}, the curated single-turn categories
of BFCL-v4 \citep{patil2025bfcl} including the category where no tool applies,
and the live categories of the same benchmark, which are queries contributed
by its users rather than written for it. Three attention designs are
represented: dense grouped-query attention at several head counts, a hybrid
stack in which only a quarter of the layers compute softmax attention at all,
and a sliding-window stack that attends locally in all but four of its layers
and carries four query heads. Failure rates span an order of magnitude, from
\rateQwenHybridGlaive{}\% to \rateLlamaOneBBfcl{}\%, which lets us separate
claims about detectors from claims about how often the models fail. The
Qwen3.5-2B run yields too few labelled failures to support a conclusion and is
reported only as a boundary case.

Each run is scored on the hardest population it supports. Where the benchmark
contains the category with no applicable tool, that is the call-expected
population, because on the remaining items the label reduces to whether a call
was emitted at all; otherwise it is every parseable call. Calls that fail to
parse because our generation budget cut them off are excluded throughout, as
they measure our protocol rather than the model.

\subsection{Head resolution}
\label{sec:exp:perhead}

\begin{table}[t]
\centering
\scriptsize
\caption{One detector family per column, scored on the hardest population
each run supports: the items where a call is genuinely required where the
benchmark has an abstention category, and every parseable call otherwise. Best
per row in bold. \emph{avg.\ spectral} averages attention over heads before
computing anything, as the prior literature does; \emph{per-head} keeps the
heads apart. \emph{surface} is a classifier on lengths alone and is not a
detector but a floor. Bold marks every detector within one seed standard
deviation of the row's leader, since the ordering inside that band changes
with the split seed; Appendix~\ref{app:spread} gives the spread.}
\label{tab:detectors}
\begin{tabular}{llcccccccc}
\toprule
Model & Data & logprob & surface & avg.\ spectral & per-head & LapEigvals & Lookback & token-role & token-level \\
\midrule
Llama-3.2-1B & Glaive & 0.576 & 0.556 & 0.535 & 0.940 & 0.942 & 0.917 & \textbf{0.957} & \textbf{0.963} \\
Llama-3.2-3B & Glaive & 0.738 & 0.609 & 0.652 & \textbf{0.930} & 0.907 & \textbf{0.911} & \textbf{0.937} & 0.897 \\
Gemma-3-1B & Glaive & 0.647 & 0.699 & 0.716 & 0.856 & 0.854 & 0.765 & \textbf{0.915} & \textbf{0.914} \\
Llama-3.2-1B & BFCL & 0.623 & 0.754 & 0.773 & 0.873 & 0.873 & 0.861 & \textbf{0.913} & 0.903 \\
Llama-3.2-3B & BFCL & 0.692 & 0.632 & 0.757 & 0.910 & 0.895 & 0.892 & \textbf{0.929} & \textbf{0.930} \\
Qwen3-1.7B & BFCL & 0.537 & 0.599 & 0.795 & 0.892 & 0.905 & 0.911 & \textbf{0.927} & 0.908 \\
MiniCPM5-2B & BFCL & \textbf{0.805} & 0.546 & 0.635 & 0.730 & 0.722 & 0.771 & 0.788 & 0.724 \\
Qwen3.5-0.8B & BFCL & \textbf{0.769} & 0.616 & 0.599 & 0.657 & 0.627 & 0.634 & 0.708 & 0.674 \\
Qwen3.5-4B & BFCL & \textbf{0.821} & 0.637 & 0.620 & 0.665 & 0.726 & 0.692 & 0.749 & 0.737 \\
Qwen3.5-2B$^\dagger$ & Glaive & 0.784 & 0.721 & 0.708 & 0.638 & 0.678 & -- & 0.746 & -- \\
Llama-3.2-1B & BFCL-live & 0.403 & 0.709 & 0.658 & \textbf{0.816} & 0.782 & 0.800 & \textbf{0.808} & \textbf{0.808} \\
MiniCPM5-2B & BFCL-live & \textbf{0.776} & 0.582 & 0.582 & 0.725 & 0.661 & 0.739 & 0.726 & 0.741 \\
\bottomrule
\end{tabular}

\end{table}

Table~\ref{tab:detectors} is the paper's central measurement, and it isolates one variable. Every row reads the
same attention tensors of the same generations under the same splits and the same seeds; the blocks differ only
in whether the heads are averaged before the graph is built. The head-averaged column sits between
\aucHeadAveragedLayerLlamaOneBGlaiveSem{} and \aucSymEigProfileLlamaThreeBGlaiveSem{} and on several runs
does not clear the surface-length confound. The per-head block reaches
\aucPerHeadAllLlamaThreeBGlaiveSem{} on Llama-3B and \aucPerHeadStaticDynLlamaThreeBBfclCall{} on the same model
on BFCL. Averaged over the \nRunsBothSpectral{} runs where both are available, keeping the heads apart is worth
\perHeadGainMean{} AUC points, with a range of \perHeadGainMin{} to \perHeadGainMax{}; the single negative entry
is the hybrid model with too few positives to interpret.

This is what Section~\ref{sec:theory} predicts, and it is worth being clear about the direction of the
inference. Proposition~\ref{prop:annihilation} says an averaged readout cannot in principle recover
between-head contrast, and Proposition~\ref{prop:jensen} says the averaged graph systematically overstates
connectivity unless heads share a Fiedler vector, which they measurably do not. Neither statement guarantees
that a failure is visible in the contrast. Table~\ref{tab:detectors} establishes that it is.

\begin{figure}[t]
\centering
\begin{minipage}[t]{0.42\linewidth}
\centering
\includegraphics[width=\linewidth]{fig_heatmap}
\end{minipage}\hfill
\begin{minipage}[t]{0.56\linewidth}
\centering
\includegraphics[width=\linewidth]{fig_ladder}
\end{minipage}
\caption{Left: where the attention signal lives. Univariate AUC of each head's spectral radius on the
generated call, Llama-3.2-1B on Glaive, direction-free. The signal is concentrated in particular heads of
particular layers; averaging over a row of this map mixes its best head with its worst. Right: the
resolution ladder of Section~\ref{sec:exp:ladder}, mean AUC over \ladNRuns{} runs. Grey bars are the two
$L\times5$ reads, blue bars keep heads apart with an increasing number of heads, and red bars have the full
per-head width but no head information.}
\label{fig:ladder}
\end{figure}

\subsection{Resolution or dimensionality?}
\label{sec:exp:ladder}

The per-head profile has $L \times H \times 5$ features and the head-averaged block $L \times 5$, so the
comparison in Table~\ref{tab:detectors} varies two things at once. A ladder of controls separates them,
all on the same stored features, the same regularised linear classifier and the same tool-grouped splits
over \ladNRuns{} runs. \emph{Metric averaging} computes every metric per head and then averages over heads,
giving $L \times 5$ features like the literature's read but with the Jensen effect of
Proposition~\ref{prop:jensen} removed: it scores \ladMetricAvg{} against \ladGraphAvg{} for the averaged
graph, a change of \ladGainMetricAvg{} that is positive on \ladGainMetricAvgPos{} of \ladNRuns{} runs.
Keeping the heads apart adds a further \ladGainPerHead{} on top of that, to \ladPerHead{}. Two controls
hold the width fixed at the per-head value. Padding the averaged-graph features with Gaussian noise to the
same width leaves them at \ladNoisePadded{}, a change of \ladNoiseEffect{} from the unpadded value, so
width without information costs rather than buys. Permuting the head axis independently for every example, which keeps every
value and destroys only which head it came from, drops the per-head profile to \ladHeadShuffled{}, a
loss of \ladGainVsShuffle{}. Finally a random subset of one, two, four and eight heads per layer scores
\ladOneHead{}, \ladTwoHeads{}, \ladFourHeads{} and \ladEightHeads{}, rising with the number of heads
kept rather than jumping at the first. The signal is carried by which heads do what, and a read that
cannot tell the heads apart cannot see it, whatever its width.

Which of the five per-head metrics carries the signal changes with the run: the spectral radius leads on
Llama-1B with \aucPerHeadLmaxLlamaOneBGlaiveSem{} against the Fiedler value's
\aucPerHeadFiedlerLlamaOneBGlaiveSem{}, and the ordering reverses on Llama-3B. Reporting the per-run best would
be selection on test, so we recommend and report the five-metric profile, which is at or near the top in every
run and costs the same single eigendecomposition.

\subsection{The access frontier}
\label{sec:exp:frontier}

\begin{table}[t]
\centering
\small
\caption{The access frontier. One representative per tier, fixed before any
result was seen: the mean log-probability, the five-metric per-head profile
and LapEigvals for attention, the token-role probe for the residual stream,
and the length classifier as the floor. Subsets as in
Table~\ref{tab:detectors}.}
\label{tab:frontier}
\resizebox{\linewidth}{!}{%
\begin{tabular}{llccccccccc}
\toprule
Model & Data & floor & logits & $\Delta$ & attention: per-head & $\Delta$ & attention: LapEigvals & $\Delta$ & residual: token-role & $\Delta$ \\
\midrule
Llama-3.2-1B & Glaive & 0.556 & 0.576 & +0.020 & 0.940 & +0.384 & 0.942 & +0.387 & 0.957 & +0.401 \\
Llama-3.2-3B & Glaive & 0.609 & 0.738 & +0.129 & 0.930 & +0.322 & 0.907 & +0.298 & 0.937 & +0.328 \\
Gemma-3-1B & Glaive & 0.699 & 0.647 & -0.052 & 0.856 & +0.157 & 0.854 & +0.155 & 0.915 & +0.217 \\
Llama-3.2-1B & BFCL & 0.754 & 0.623 & -0.131 & 0.873 & +0.119 & 0.873 & +0.120 & 0.913 & +0.160 \\
Llama-3.2-3B & BFCL & 0.632 & 0.692 & +0.059 & 0.910 & +0.277 & 0.895 & +0.263 & 0.929 & +0.296 \\
Qwen3-1.7B & BFCL & 0.599 & 0.537 & -0.062 & 0.892 & +0.294 & 0.905 & +0.306 & 0.927 & +0.328 \\
MiniCPM5-2B & BFCL & 0.546 & 0.805 & +0.259 & 0.730 & +0.184 & 0.722 & +0.176 & 0.788 & +0.241 \\
Qwen3.5-0.8B & BFCL & 0.616 & 0.769 & +0.152 & 0.657 & +0.040 & 0.627 & +0.011 & 0.708 & +0.091 \\
Qwen3.5-4B & BFCL & 0.637 & 0.821 & +0.184 & 0.665 & +0.028 & 0.726 & +0.089 & 0.749 & +0.112 \\
Qwen3.5-2B$^\dagger$ & Glaive & 0.721 & 0.784 & +0.064 & 0.638 & -0.083 & 0.678 & -0.042 & 0.746 & +0.025 \\
Llama-3.2-1B & BFCL-live & 0.709 & 0.403 & -0.306 & 0.816 & +0.107 & 0.782 & +0.074 & 0.808 & +0.099 \\
MiniCPM5-2B & BFCL-live & 0.582 & 0.776 & +0.194 & 0.725 & +0.143 & 0.661 & +0.078 & 0.726 & +0.143 \\
\bottomrule
\end{tabular}
}

\end{table}

Table~\ref{tab:frontier} is the deployment result, and it has two regimes. On the Llama, Qwen3 and Gemma runs
the ordering of the internal tiers is stable: the residual stream is best, attention-only trails it by a few
points, and both are far above the confound floor, while the token log-probability reaches
\aucLogprobLlamaThreeBBfclCall{} on Llama-3B and \aucLogprobQwenDenseBfclCall{} on Qwen3-1.7B, which is chance,
on a run where the same model's residual stream supports \aucHiddenLRQwenDenseBfclCall{}. A guardrail built on
output confidence would pass exactly the calls this paper is about on those models. On the MiniCPM5 and Qwen3.5
runs the picture inverts: the log-probability is the best single detector and no trained judge reaches it.
Section~\ref{sec:exp:budget} examines that division. Within the internal tiers the token-role probe leads
the per-head profile on eleven of the twelve runs; on the one where attention leads, the difference is
within seed noise.
The weakness of confidence on the first group does not track model quality: Qwen3-1.7B abstains on the
no-applicable-tool category almost as often as any model in the sweep (Section~\ref{sec:exp:arch}), and its
confidence is the least informative we measured.

Within the attention tier, LapEigvals and our per-head profile trade places by run and stay within noise of one
another. We do not claim to beat it. What the per-head profile offers instead is compactness, $L\times H\times 5$
features against $L\times H\times k$ with $k$ up to $100$, and a readout with a graph-theoretic meaning, which is
what makes Section~\ref{sec:theory} possible. A smaller probe should also need fewer labelled failures, and
Figure~\ref{fig:labeleff} tests that by thinning the positives available to the training folds on
\leNRuns{} runs while scoring on every example. The two attention judges are tied at full budget
(\lePerHeadFull{} against \leLapEigFull{}) and at the smallest budget, a handful of failures per fold
(\lePerHeadFive{} against \leLapEigFive{}); in between, with a tenth or a fifth of the labels, the per-head
profile leads by \leDiffTen{} and \leDiffTwenty{}, on \leWinsTen{} and \leWinsTwenty{} of \leNRuns{} runs.
The edge is modest and not uniform across runs, so we state it as parity with a small advantage in the
low-label regime, not as a general label-efficiency result. The residual probe leads both throughout,
\leHiddenTen{} at a tenth of the labels against \lePerHeadTen{}, and it loses less than either as labels
shrink.

\begin{figure}[t]
\centering
\includegraphics[width=0.5\linewidth]{fig_labeleff}
\caption{Label efficiency of the trained judges. Positives are thinned in the training folds only, every
example is scored, and the curve is the mean over \leNRuns{} runs of the Llama, Gemma and Qwen3 families.
The two attention judges are tied at both ends and the per-head profile leads by a few points in between;
the residual probe leads throughout.}
\label{fig:labeleff}
\end{figure}

\subsection{Decorrelation of the two attention families}
\label{sec:exp:decorrelate}

Corollary~\ref{cor:decorrelate} makes three predictions from Remark~\ref{prop:lapeig}. Two hold and one
fails, and the failure carries the more useful lesson.

The first prediction, that the two families' scores are not interchangeable, holds. Over \decorrNRuns{} runs the
Spearman correlation between the per-head spectral score and the LapEigvals score averages \decorrRhoMean{},
ranging from \decorrRhoMin{} to \decorrRhoMax{}. They agree substantially, as two competent detectors of the same
events must, but they are far from being the same function of the forward pass.

The second prediction, that each retains signal once the other is removed, holds strongly. Regressing the
per-head score on the LapEigvals score and scoring the residual gives \decorrResidPerHead{} AUC on average, and
the reverse gives \decorrResidLapEig{}. Both are far above chance, so neither family is a noisy copy of the
other: an in-degree sequence and a symmetrised spectrum carry substantially independent evidence about the same
call, as Proposition~\ref{prop:lapeig} implies they should.

The third prediction, that a combination should therefore exceed both, fails. Concatenating the two feature sets
changes AUC by \decorrUnionGain{} on average and improves on the better of the two in only
\decorrUnionPositiveRuns{} of \decorrNRuns{} runs, with a best case of \decorrUnionGainMax{}. For feature
concatenation a variance account is at least available: the union carries thousands of features against a few
hundred positives. The two-score combinations of Section~\ref{sec:exp:agreement} have no such excuse and fail
too, so that account cannot be the whole story. The distinction matters for practice. Earlier work reported that
combining internal signals does not help and read that as evidence that the signals are redundant. The residual
analysis shows they are not redundant in the sense of being the same function of the forward pass. Whether they are redundant for the
decision is a different question, and Section~\ref{sec:exp:agreement} answers it in the negative direction.

Because this is a negative result that happens to favour our own framing, the union is given the more generous
of two protocols: its LapEigvals width is selected on validation rather than fixed. Fixing it instead gives
\decorrUnionGainFixed{}, so the conclusion does not depend on that choice.

\subsection{Combining the two judges}
\label{sec:exp:agreement}

Section~\ref{sec:exp:decorrelate} shows that the residual-stream and attention
judges carry substantially independent evidence, which makes it natural to ask
whether running both is worth the cost. We tried three ways of combining them
and none of them is.

Concatenating their features changes AUC by \decorrUnionGain{}. Stacking their
outputs with a combiner fitted on out-of-fold training scores, which is the
standard remedy for a small validation fold and gives the combiner several
times more labelled failures to work with, changes AUC by \agrStackGain{} and
improves on the better single judge in \agrStackPositive{} of \agrNRuns{}
runs. Putting both scores on one scale first, by mapping each through the
empirical distribution of its out-of-fold training scores, changes AUC by
\agrStackRankGain{} and helps in \agrStackRankPositive{} of \agrNRuns{}
runs, with a best case of \agrStackRankGainMax{}; every fitted combiner
gave both judges positive weight. Part of the remaining deficit is the
pooling penalty of Section~\ref{sec:exp:stats}, which a combination pays
twice, so we also score the simplest combination inside each fold, where no
such penalty exists: on
Llama-1B on Glaive the average of the two judges' scores reaches
\poolWithinMeanLlamaOneBGlaive{} against \poolWithinPerHeadLlamaOneBGlaive{}
for the better single judge in the same folds, and on Llama-3B
\poolWithinMeanLlamaThreeBGlaive{} against \poolWithinHiddenLlamaThreeBGlaive{}.
The combination matches the better judge and does not exceed it.

The third is the one an operator would reach for without any fitting: set a
threshold for each judge on training data and flag a call only when both fire.
It appears to work, and the appearance is an artefact of the operating point.
The rule fires on fewer calls than either judge alone, so it runs at lower
recall, and lower recall buys precision for nothing. Compared against the
stronger single judge \emph{at the recall the rule actually achieves}, the
agreement rule changes precision by \agrMatchedGain{} on average and helps in
\agrMatchedPositive{} of \agrNRuns{} runs. Moving one judge's threshold does
the same work.

We therefore report a negative result without a mechanism we can defend. The
two families are far from interchangeable, and yet nothing we tried converts
that into a better detector, within folds or pooled. A shortage of labels is
not the reason: the combiners here had hundreds of out-of-fold positives, and
Section~\ref{sec:exp:budget} shows a far larger probe surviving a tenfold
thinning of its labels. The likelier reading is that the two judges disagree
mostly on examples both already score correctly, and agree on the ones both
get wrong. The recommendation that follows is to pick the judge the access
tier allows and spend the labelling effort on its threshold rather than on a
second judge.

\paragraph{Precision is bounded by how often the model fails.}
One property of these operating points matters more for deployment than any
difference between detectors. Across our runs the failure rate ranges from
\agrBaseMin{}\% to \agrBaseMax{}\%, and precision tracks it: on the models
that fail most, a judge reaches precision in the high nineties at usable
recall, and on the model that fails least, no judge we measured exceeds a
fifth. A guardrail is therefore most deployable exactly where the model is
worst, and on a well-behaved model the false positives it raises may cost more
than the failures it catches. That is an argument for measuring the failure
rate of the deployed model before choosing a guardrail, and it is the same
measurement Section~\ref{sec:exp:budget} already recommends for a different
reason.

\subsection{Architecture and abstention}
\label{sec:exp:arch}

Attention design changes what the attention tier can read, and the effect is smaller than the earlier literature
suggests. On the sliding-window model, which has four heads and one key-value head and computes local attention
in all but four of its layers, the per-head profile still reaches \aucPerHeadAllGemmaGlaiveSem{} against a
confound floor of \aucSurfaceConfoundGemmaGlaiveSem{}. A narrow, sparse attention stack therefore does not
disable spectral detection. The hybrid families, in which only a quarter of the layers carry softmax attention,
are where the tier weakens, and there the attention-only detectors fall to the neighbourhood of the residual
probes rather than below the confound floor.

The clearest architectural effect we find is not about detectors but about the models. On the benchmark category
where no available tool applies and the correct behaviour is to answer without calling anything, the two Llama
models fabricate a call almost every time: \modeLlamaOneBBfclOvertrigger{} of
\nIrrelevanceItems{} items at 1B and \modeLlamaThreeBBfclOvertrigger{} at 3B, with scale making it worse rather
than better. Every other family we tested abstains most of the time, and abstains at every size we tested it:
Qwen3-1.7B on \modeQwenDenseBfclValidnocall{} items, Qwen3.5-4B on \modeQwenFourBBfclValidnocall{}, MiniCPM5-2B
on \modeMiniCpmBfclValidnocall{}, and the smallest model in the study, Qwen3.5-0.8B, on
\modeQwenSmallBfclValidnocall{}. We inspected these decisions by hand: the Llama outputs are well-formed calls,
and the others are prose answers to the question that was asked.

Abstention is therefore a property of the training recipe and not of scale, and an earlier version of this study,
which had only the Llama family in it, drew the opposite and wrong conclusion that small models cannot abstain.
The practical consequence is unchanged in form but narrower in scope: for a model family that does not abstain,
the capability has to be supplied from outside the model, and a guardrail of the kind measured here is the only
option. For a family that does abstain, the guardrail is guarding against a rarer event, which is exactly the
regime that Section~\ref{sec:exp:budget} shows makes a trained judge hard to fit.

\subsection{Confounds and transfer}
\label{sec:exp:confound}

Every table reports the confound floor beside the detectors, and the floor is not always low: on Llama-1B on
BFCL a classifier built from lengths alone reaches \aucSurfaceConfoundLlamaOneBBfclAll{} on the full set,
because format failures are long. That is why the semantic and call-expected subsets exist, and on them the
floor falls to \aucSurfaceConfoundLlamaOneBBfclCall{} while the detectors hold.

Transfer separates a detector that learned routing from one that learned a corpus. Trained on Glaive and applied
unchanged to BFCL on the same model, the residual probe retains
\transferHiddenLRLlamaThreeBGlaiveToLlamaThreeBBfcl{} on Llama-3B and the per-head profile
\transferPerHeadAllLlamaThreeBGlaiveToLlamaThreeBBfcl{}, against in-domain values of
\aucHiddenLRLlamaThreeBBfclSem{} and \aucPerHeadAllLlamaThreeBBfclSem{}. The length confound does not transfer at
all: it falls to \transferSurfaceConfoundLlamaThreeBGlaiveToLlamaThreeBBfcl{}, below chance, on the same
comparison. A feature that was carrying corpus-specific length regularities behaves exactly as one would expect
under distribution shift, and the detectors do not.

\begin{figure}[t]
\centering
\begin{minipage}[t]{0.5\linewidth}
\centering
\includegraphics[width=\linewidth]{fig_family}
\end{minipage}\hfill
\begin{minipage}[t]{0.46\linewidth}
\centering
\includegraphics[width=\linewidth]{fig_thinning}
\end{minipage}
\caption{Left: the advantage of the best residual probe over mean log-probability on every run, scored on
every parseable call. The two groups of families do not overlap. Right: the probe's AUC on Llama-3.2-1B on
BFCL as its labelled failures are thinned, against the confidence baseline on the same data. A tenfold
reduction does not bring the probe near the baseline, so the split on the left is not a shortage of labels.}
\label{fig:family}
\end{figure}

\subsection{Calls with a history}
\label{sec:exp:multiturn}

Every result so far concerns a single call made from a fresh prompt, and an agent does not work that way: its
calls are made with a conversation behind them, and a wrong call can be inherited from an earlier one. We
therefore evaluate on the multi-turn categories of BFCL-v4 under a teacher-forced protocol. A conversation is
replayed with the ground-truth calls of the two preceding turns in the prompt, the model is asked for the calls
of the current turn, and the label is whether those calls match the ground truth for that turn. This keeps the
label unambiguous without running the benchmark's simulated back ends, at the price that the model never sees
its own earlier mistakes. To restore that condition deliberately, half of the records with a history carry one
preceding turn whose calls were replaced by a wrong call drawn from another conversation, presented as having
succeeded, which we call the \emph{cascade} condition. Two models are run, one from each side of the family
split of Section~\ref{sec:exp:budget}: Llama-3.2-1B on \mtNLlamaOneB{} turns and MiniCPM5-2B on
\mtNMiniCpm{}. Prompts run to several thousand tokens, so the attention features are reduced layer by layer
inside a forward hook rather than materialised for the whole stack (Appendix~\ref{app:protocol}). Because a
weak model echoes tool schemas into argument-less calls here, labels on these runs count arguments the schema
does not define as a failure (Appendix~\ref{app:validity}).

\paragraph{The models.}
Multi-turn calling is far harder for these models than the single-turn categories. MiniCPM5 fails on
\mtFailTurnZeroMiniCpm{}\% of opening turns and \mtFailAllMiniCpm{}\% of all turns; Llama-1B fails on
\mtFailTurnZeroLlamaOneB{}\% of opening turns, \mtTruncPctLlamaOneB{}\% of its generations run past the budget,
and the whole run leaves \mtNegLlamaOneB{} well-formed correct calls, too few to train or score a judge on. We
report Llama-1B as a boundary case, the multi-turn counterpart of the run in Table~\ref{tab:runs} that fails
too rarely: a judge needs both kinds of example, and a model that almost never succeeds supplies no more of
them than one that almost never fails. The cascade condition asks whether a wrong call in the history makes the
next call worse, and on neither model does it. On turns with a history MiniCPM5 fails on
\mtFailCleanMiniCpm{}\% after a clean history and \mtFailCorruptMiniCpm{}\% after a corrupted one (Fisher
$p=\mtFisherPMiniCpm{}$ on \mtNCleanMiniCpm{} and \mtNCorruptMiniCpm{} turns), and Llama-1B on
\mtFailCleanLlamaOneB{}\% against \mtFailCorruptLlamaOneB{}\% ($p=\mtFisherPLlamaOneB{}$). A single wrong
call that the history reports as successful does not propagate; whether a wrong \emph{result} would is a
question this protocol cannot ask, and we return to it in Section~\ref{sec:disc}.

\paragraph{The judges.}
Detection survives the history. On MiniCPM5, with tools held out as before, the token-role probe reaches
\mtAucHiddenMiniCpm{}, the per-head profile \mtAucPerHeadMiniCpm{} and LapEigvals \mtAucLapEigMiniCpm{}, against
a length floor of \mtAucSurfaceMiniCpm{}; the probe scores turns with a history at \mtAucHiddenHistMiniCpm{}
and opening turns at \mtAucHiddenTurnZeroMiniCpm{}, so the history neither helps nor hurts it. The family
split of Section~\ref{sec:exp:budget} reappears intact: mean log-probability reaches \mtAucLogprobMiniCpm{} on
this model and no trained judge exceeds it, with a history behind the call as without one. On Llama-1B every
judge sits near chance on the \mtNegLlamaOneB{} negatives available, LapEigvals highest at
\mtAucLapEigLlamaOneB{}, and we draw no conclusion from that run.

\paragraph{The judge under cascade.}
The question a deployment cares about is whether a judge trained on well-behaved histories still works once the
history contains a failure it has never seen. We train every judge on clean-history items only and score it, on
held-out tools, separately on clean and on corrupted items. On MiniCPM5 the token-role probe scores
\mtAucHiddenCleanOnCleanMiniCpm{} on clean and \mtAucHiddenCorruptOnCleanMiniCpm{} on corrupted items, the
per-head profile \mtAucPerHeadCleanOnCleanMiniCpm{} against \mtAucPerHeadCorruptOnCleanMiniCpm{}, and LapEigvals
\mtAucLapEigCleanOnCleanMiniCpm{} against \mtAucLapEigCorruptOnCleanMiniCpm{}. A corrupted history does not
change what a wrong call looks like from inside the model, at least not in a way that a judge trained without
ever seeing one would miss; training on both conditions gives \mtAucHiddenCorruptOnAllMiniCpm{} on the
corrupted items, the same figure. For an operator this is the useful half of a null result: the calibration
set for a judge does not need to contain failed trajectories.

\subsection{When a trained judge beats output confidence}
\label{sec:exp:budget}

Whether a trained internal judge is worth building at all is the first
question an operator asks, and our runs do not give it one answer. They
divide cleanly by model family (Figure~\ref{fig:family}). On the Llama, Qwen3 and Gemma runs a probe on
the residual stream leads mean log-probability by a wide margin. On the
MiniCPM5 and Qwen3.5 runs the margin nearly vanishes, and on the harder of
our two evaluation subsets it reverses and confidence wins.

The division is complete rather than a tendency, and we are careful about
the unit over which that is claimed. Over \famNRunsSem{} runs scored on
every parseable call, the probe's advantage averages \famGapRecentSem{} on
the recent families against \famGapEarlierSem{} on the others, and the two
groups do not overlap: the largest advantage among the recent families,
\famRecentMaxSem{}, falls below the smallest among the others,
\famEarlierMinSem{}. Runs are not independent observations, however, since
one checkpoint is evaluated on up to three benchmarks, so the unit for any
test is the checkpoint, of which there are \famNCkptSem{},
\famNCkptRecentSem{} recent and \famNCkptEarlierSem{} earlier. They too
separate completely, which is the smallest attainable two-sided rank-test
value at that size, $p=\famCkptMannWhitneySem{}$. At the level of the
family, which is what the claim is about, there are five observations and
no test has power; we report the family means instead, Qwen3
\famGapQwenThreeSem{}, Gemma-3 \famGapGemmaThreeSem{} and Llama-3.2
\famGapLlamaThreeTwoSem{} against MiniCPM5 \famGapMiniCPMFiveSem{} and Qwen3.5
\famGapQwenThreeFiveSem{}, and let the reader judge the gap. Restricting to the
items where a call is genuinely required, which removes the abstention
decision and is the population this paper is about, the recent families
average \famGapRecentCall{}: confidence beats every trained detector we
measured on \famConfWinsCall{} of \famNRecentCall{} such runs. That subset
has only \famNCkptCall{} checkpoints, so its separation is descriptive. We
report both subsets because the sign of the effect for the recent families
depends on which is used, while the separation between families does not.

The obvious explanation is not the right one, and we report the test because
we expected it to succeed. The recent families also fail far less often, and
a supervised probe needs labelled failures while confidence needs none, so
their probes might simply be starved. Two checks rule that out. First,
thinning the labelled failures available to the Llama-1B run on BFCL from
\curveBaseLlamaOneBBfclFullPos{} to \curveBaseLlamaOneBBfclLowPos{}, a
tenfold reduction, costs the probe only \curveBaseLlamaOneBBfclDrop{} AUC,
from \curveBaseLlamaOneBBfclFull{} to \curveBaseLlamaOneBBfclLow{}, which
still leaves it far above the \confBaseLlamaOneBBfcl{} that confidence reaches
on the same data. A probe trained on twenty-five failures wins comfortably, so
scarcity of labels is not what holds the other families back. Second, the
apparent relationship between failures and advantage is an artefact of
grouping. Across all runs the rank correlation between the number of labelled
failures and the probe's advantage is \famSpearmanSem{} and does not reach
significance at $p=\famSpearmanPSem{}$, and computed within each family
separately it disappears, at \famWithinRecentSem{} and
\famWithinEarlierSem{}. One run settles it alone: Qwen3.5-0.8B supplies
\posQwenSmallBfclCount{} labelled failures, more than most runs where the
probe wins, and confidence still leads there.

We therefore attribute the difference to the model family. The recent
families we tested express uncertainty that tracks their own tool-calling
errors, so their output distribution is already informative and a trained
judge adds little; the older families do not, and there a trained judge is
the only thing that works. With five families we cannot say which property of
training produces this, and we do not claim a trend over time. Three
confounds remain open and we name them rather than argue them away. The two
recent families both emit their calls in an XML dialect while the others emit
JSON, so the token probabilities being summarised are taken over different
surface forms; a control that forces one family into the other's format
would settle this and we have not run it. One of the two recent families has
a hybrid attention stack, but the other, MiniCPM5, is a dense stack on which
confidence wins just as clearly, so attention design is not the explanation.
Finally, a model that has seen the benchmark in training will fail less and
know when it fails, which is what we observe; the BFCL-live categories are
contributed by users after the curated set and are the less likely to have
been seen, and on them MiniCPM5 shows the same behaviour, which weakens the
contamination account without excluding it. The practical
consequence is nevertheless firm, and it is not the rule we set out to write:
before building an internal guardrail, measure output confidence on a few
hundred labelled calls from the model that will actually be deployed, because
whether internals are needed is a property of that model and cannot be
inferred from its size, its attention design, or its release date.

\subsection{Does the comparison depend on how confidence is summarised?}
\label{sec:exp:confvariants}

Comparing a trained judge against ``output confidence'' presumes a way of
summarising the output distribution, and the mean token log-probability is
only one. If a better summary closed the gap on the families where confidence
loses, the split of Section~\ref{sec:exp:budget} would be an artefact of that
choice rather than a property of the models. We therefore record
\CvLlamaBfclNSummaries{} summaries per generated call, covering the mean,
minimum, total and mean-lowest-five log-probability, perplexity, the fraction
of low-probability tokens, and the mean, maximum and final predictive
entropy, and score each as its own detector under the identical protocol with
its direction fixed on training folds.

The choice matters, and it matters in the direction that weakens our own
earlier framing. On Llama-3.2-1B on BFCL the mean log-probability reaches
\CvLlamaBfclMeanLp{} while the best summary that is not confounded with
length, the mean of the five least likely tokens, reaches \CvLlamaBfclBest{}.
A single improbable token is what carries the evidence, and averaging over the
call dilutes it. The probe's advantage therefore falls from
\CvLlamaBfclGapMean{} to \CvLlamaBfclGapBest{}. It does not disappear: the
best available confidence summary still trails the probe's
\CvLlamaBfclProbe{} by a wide margin. One summary must be set aside. The total
log-probability of the span scores \CvLlamaBfclSumLp{}, which is close to the
\CvLlamaBfclSurface{} of the length baseline on the same run, and since it
grows with length it is largely measuring the same thing; we report it but do
not treat it as a confidence result.

On the family where confidence wins, the choice of summary makes almost no
difference. On MiniCPM5-2B the best summary reaches \CvMiniCpmBfclBest{}
against \CvMiniCpmBfclMeanLp{} for the mean log-probability, and the probe's
deficit moves from \CvMiniCpmBfclGapMean{} to \CvMiniCpmBfclGapBest{}. The
summaries also agree with one another there, most of them falling within a
couple of points, whereas on Llama they range over more than seventeen. That
is a second and independent way of stating the split: on one family the
model's expressions of uncertainty agree with each other and with its errors,
and on the other they neither agree nor track the errors.

The consequence for the rest of the paper is a caveat rather than a
correction. The cross-run comparison of Section~\ref{sec:exp:budget} uses the
mean log-probability, which we now know understates confidence on the
families where it loses; the separation between families survives the fairer
comparison, but it is smaller than a mean-log-probability comparison alone
would suggest.

\subsection{Statistical discipline}
\label{sec:exp:stats}

\paragraph{Pooled scores are conservative.}
Every AUC in this paper pools the out-of-fold scores of five classifiers, one per fold, each calibrated on
its own training tools. A deployed judge is one classifier with one calibration, so pooling penalises a
calibration drift between folds that a deployment never sees. We measured the cost on \poolNRuns{} runs by
scoring the residual probe both ways: within folds it averages \poolWithinHiddenLlamaOneBGlaive{} on
Llama-1B on Glaive against \poolPooledHiddenLlamaOneBGlaive{} pooled, and the cost averages
\poolCostHiddenMean{} AUC across the runs. We keep the pooled protocol because per-fold AUCs on this data
are computed on a handful of positives and are unstable, but the reader should take every number in
Tables~\ref{tab:detectors} and~\ref{tab:frontier} as a floor on what a single calibrated judge achieves.
The same effect is much larger for any score that is only meaningful relative to its fold, such as a rank:
the within-fold rank average of the residual and attention judges scores
\poolWithinRankMeanLlamaOneBGlaive{} inside folds and \poolPooledRankMeanLlamaOneBGlaive{} pooled, which is
why Section~\ref{sec:exp:agreement} reads its combinations within folds.

\paragraph{Ties.}
Differences within the top group of detectors are smaller than the variation between split seeds. On Llama-1B a
paired bootstrap over the pooled scores puts the residual probe, LapEigvals and the per-head profile within a
three-way tie whose ordering changes with the seed. We therefore state the frontier between tiers, which is
stable and large, and refrain from ranking detectors inside a tier, which is not.
